\chapter{Geltungsbereich und elementare Folgerungen}

Dieser Leitfaden sammelt zentrale Definitionen und Sätze für beliebige
Halbgruppen $S$. Vorausgesetzt werden nur eine überall definierte innere
zweistellige Verknüpfung und Assoziativität: keine Endlichkeit, Einheit,
Null, Kommutativität, Kürzbarkeit, Idempotenz oder Regularität.

Spezialsätze über vorgegebene idempotente oder reguläre Elemente bleiben
ausgenommen. Allgemeine Konstruktionen und vollständige Fallunterscheidungen
sind enthalten. Dabei können Monoide oder Gruppen als Ergebnisse auftreten;
neutrale oder inverse Elemente in $S$ werden nicht vorausgesetzt.

Der Band enthält weder Beweise noch Beispiele und beansprucht keine
Vollständigkeit. Halbgruppen und Unterhalbgruppen dürfen leer sein;
Ideale und Untergruppen sind nichtleer. Als mengentheoretische
Grundlage gilt ZFC einschließlich des Auswahlaxioms. Das Auswahlaxiom
ist damit im gesamten Leitfaden zugelassen und stellt keine zusätzliche
Eigenschaft der betrachteten Halbgruppen dar.

\section{Grundbegriff und Notation}

\begin{definition}[Halbgruppe und positive Produkte]
Eine \emph{Halbgruppe} ist eine Menge $S$ mit einer Abbildung
\[
 \mu:S\times S\longrightarrow S,\qquad (a,b)\longmapsto ab,
\]
für die $(ab)c=a(bc)$ für alle $a,b,c\in S$ gilt.
Die Abbildungseigenschaft enthält die \emph{Abgeschlossenheit}.
Im ganzen Leitfaden bezeichnen $S,T,U$ beliebige Halbgruppen.
Es sei $\mathbb N_{>0}=\{1,2,\ldots\}$.
Für $a\in S$ werden die \emph{positiven Potenzen} durch
$a^1=a$ und $a^{n+1}=a^na$ definiert.
Ein leeres Produkt und eine Potenz $a^0$ werden in $S$ nicht vorausgesetzt.
\end{definition}

\begin{definition}[Äußere Einheit und Kontextnotation]
\label{b28guide:def:unit}
Es sei $S^1=S$, wenn $S$ ein zweiseitiges neutrales Element besitzt;
dieses wird dann mit $1$ bezeichnet.
Andernfalls bezeichnet $S^1=S\mathbin{\dot\cup}\{1\}$ die Erweiterung
mit $1x=x1=x$ für alle $x\in S^1$.
Damit dürfen in Ausdrücken $xay$ die linken oder rechten Faktoren
durch $1$ entfallen. Eine Einheit in $S$ wird dadurch nicht verlangt.
\end{definition}

\section{Klammerungen und Potenzen}

\begin{theorem}[Allgemeines Assoziativgesetz]
\label{b28guide:thm:bracketing}
Für jedes $n\geq1$ und jede Folge $a_1,\ldots,a_n\in S$
haben alle vollständigen binären Klammerungen unter Beibehaltung
der Reihenfolge denselben Wert. Dieser wird mit $a_1\cdots a_n$
bezeichnet. Für $1\leq k<n$ gilt
\[
 a_1\cdots a_n=(a_1\cdots a_k)(a_{k+1}\cdots a_n).
\]
Umgekehrt ist eine überall definierte innere zweistellige Verknüpfung,
deren Produkte diese Klammerungsunabhängigkeit besitzen, assoziativ.
\cite{projectbracketing,harju}
\end{theorem}

\begin{theorem}[Positive Potenzgesetze]
Für $a\in S$ und $m,n\geq1$ gelten
\[
 a^{m+n}=a^ma^n,\qquad (a^m)^n=a^{mn},
 \qquad a^ma^n=a^na^m.
\]
Insbesondere ist $\langle a\rangle=\{a^n:n\geq1\}$ eine
kommutative Unterhalbgruppe von $S$; die gesamte Halbgruppe
muß nicht kommutativ sein. \cite{harju,gould}
\end{theorem}

\begin{theorem}[Vollständige Alternative für eine Potenzfolge]
\label{b28guide:thm:monogenic}
Für jedes $a\in S$ tritt genau einer der folgenden Fälle ein:
\begin{enumerate}
\item Alle positiven Potenzen von $a$ sind verschieden. Die Abbildung
$n\mapsto a^n$ ist ein Isomorphismus
$(\mathbb N_{>0},+)\to\langle a\rangle$.
\item Es gibt eindeutig bestimmte Zahlen $r,p\geq1$, so daß für alle $m,n\geq1$ gilt:
\[
 a^m=a^n
 \iff
 \bigl(m=n\bigr)\ \text{oder}\
 \bigl(m,n\geq r\ \text{und}\ m\equiv n\pmod p\bigr)
.
\]
Dann besteht $\langle a\rangle$ aus genau $r+p-1$ Elementen.
Die Zahlen $r$ und $p$ heißen \emph{Index} und \emph{Periode} von $a$.
\end{enumerate}
\cite{harju,gould}
\end{theorem}
